%% ---------------------------------------------------------------------
\section{Routh-Hurwitz para diseñar}
\label{sec:es-diseno}
%% ---------------------------------------------------------------------

La gran ventaja de Routh-Hurwitz sobre \cmd{roots} es que trabaja con los coeficientes, y los coeficientes pueden contener parámetros del controlador. La tabla responde entonces la pregunta de diseño más básica: \emph{¿para qué valores de la ganancia el lazo es estable?}

\subsection{Rango de ganancia, ganancia crítica y frecuencia de oscilación}

\begin{ejemplo}[label=ej:es-rangokc]{Rango de $\Kc$ de un lazo de proceso}
En el lazo de la \cref{fig:es-tres}, un controlador proporcional $\Kc$ actúa sobre una válvula, un proceso y un sensor de primer orden. Halle el rango de $\Kc$ para el cual el lazo es estable, la ganancia crítica y la frecuencia a la que oscila.
\solucion \textbf{1. Ecuación característica.}
\[
  1+\Kc\,\frac{0.016}{3s+1}\cdot\frac{50}{30s+1}\cdot\frac{1}{10s+1}=0
  \;\Longrightarrow\;
  900s^3+420s^2+43s+(1+0.8\Kc)=0 .
\]
\textbf{2. Tabla.}
\[
\begin{array}{c|cc}
  s^3 & 900 & 43\\
  s^2 & 420 & 1+0.8\Kc\\
  s^1 & \dfrac{17\,160-720\Kc}{420} & \\[6pt]
  s^0 & 1+0.8\Kc
\end{array}
\]
\textbf{3. Condiciones.} Fila $s^0$: $1+0.8\Kc>0$, o sea $\Kc>-1.25$. Fila $s^1$: $17\,160-720\Kc>0$, o sea $\Kc<23.83$. Entonces
\[
  \alerta{-1.25<\Kc<23.83}.
\]
\textbf{4. Ganancia y frecuencia críticas.} En $\Kc=\Kcr=23.83$ la fila $s^1$ se anula: es una fila de ceros, y el auxiliar es la fila $s^2$,
\[
  420s^2+(1+0.8\cdot23.83)=420s^2+20.07=0\;\Longrightarrow\;s=\pm j\,0.2186 .
\]
El lazo oscila sin amortiguarse con $\wcr=0.2186$\,rad/s, es decir, con periodo $\Pcr=2\pi/\wcr=28.7$\,s (\cref{fig:es-rangokc}).
\end{ejemplo}

\begin{figure}[ht]
\centering
\includegraphics[width=0.92\linewidth]{cap05/lazo_tres}
\caption{Lazo de proceso del ejemplo~\ref{ej:es-rangokc}.}
\label{fig:es-tres}
\end{figure}

\begin{figure}[ht]
\centering
\includegraphics[width=0.85\linewidth]{cap05/f_rango_kc}
\caption{Ejemplo~\ref{ej:es-rangokc}: por debajo de $\Kcr$ la oscilación se amortigua, en $\Kcr$ se sostiene con periodo $\Pcr=28.7$\,s y por encima crece.}
\label{fig:es-rangokc}
\end{figure}

\begin{ideaclave}
\textbf{Ganancia crítica: tres caras del mismo número.}\enspace $\Kcr$ y $\wcr$ salen de la fila que se anula y del polinomio auxiliar. Son los mismos números que (i) el método de oscilación sostenida de Ziegler y Nichols busca experimentalmente en la planta (capítulo~8)~\citep{ZiNi42}, y (ii) el margen de ganancia y la frecuencia de cruce de fase que se leen en el diagrama de Bode (capítulo~10): para este lazo, \cmd{margin} da $\GM=23.83$ y $\wgm=0.2186$\,rad/s.
\end{ideaclave}

\codigorepo{cap05_estabilidad/C5_Ej4_RangoK.m}{rangokc}{ejemplo~\ref{ej:es-rangokc}: rango de $\Kc$, $\Kcr$, $\wcr$ y la comparación con \cmd{margin}}

\begin{cuidado}[Las ganancias negativas]
El rango incluye $-1.25<\Kc<0$. Con $\Kc$ negativo el controlador empuja en la dirección equivocada; el lazo resulta estable mientras el término independiente siga positivo, pero el error en estado estacionario es peor que en lazo abierto. Que un valor sea estable no lo hace útil: en la práctica se usa $0<\Kc<\Kcr$, y bastante por debajo de $\Kcr$.
\end{cuidado}

\subsection{Estabilizar una planta inestable}

Cuando la planta es inestable, la ganancia ya no tiene una cota superior sola: también aparece una cota \emph{inferior}, porque con poca ganancia el controlador no alcanza a contener el polo inestable.

\begin{ejemplo}[label=ej:es-inestable]{Una ventana estrecha}
La planta $\Gp(s)=\dfrac{1}{(s-1)(s^2+2s+2)}$ se controla con el compensador $\Gc(s)=\dfrac{K(s+2)}{s+3}$, con retroalimentación unitaria (\cref{fig:es-inestable-lazo}). Halle el rango de $K$ que estabiliza el lazo.
\solucion La ecuación característica es $(s+3)(s-1)(s^2+2s+2)+K(s+2)=0$:
\[
  s^4+4s^3+3s^2+(K-2)\,s+(2K-6)=0 .
\]
\[
\begin{array}{c|ccc}
  s^4 & 1 & 3 & 2K-6\\
  s^3 & 4 & K-2 & \\
  s^2 & \dfrac{14-K}{4} & 2K-6\\[6pt]
  s^1 & \dfrac{68-16K-K^2}{14-K} & \\[6pt]
  s^0 & 2K-6
\end{array}
\]
Las condiciones son $K<14$ (fila $s^2$), $2K-6>0$ (fila $s^0$) y $68-16K-K^2>0$ (fila $s^1$, con $14-K>0$). Las raíces de $K^2+16K-68$ son $-8\pm2\sqrt{33}$, es decir $3.489$ y $-19.49$. Juntando:
\[
  \alerta{3<K<3.489}.
\]
Es un rango muy estrecho (\cref{fig:es-inestable}): un error de pocos puntos porcentuales en la ganancia, o en el modelo, desestabiliza el lazo. Este compensador estabiliza, pero no es robusto; dentro del rango, la mejor ganancia es $K\approx3.07$, con los polos más lentos en $-0.28$.
\end{ejemplo}

\begin{figure}[ht]
\centering
\includegraphics[width=0.78\linewidth]{cap05/lazo_inestable}
\caption{Ejemplo~\ref{ej:es-inestable}: compensador para una planta con un polo en $+1$.}
\label{fig:es-inestable-lazo}
\end{figure}

\begin{figure}[ht]
\centering
\includegraphics[width=0.85\linewidth]{cap05/f_inestable}
\caption{Ejemplo~\ref{ej:es-inestable}: parte real del polo más a la derecha en función de $K$. Por debajo de $3$ queda un polo real inestable; por encima de $3.489$ cruzan el eje dos polos complejos.}
\label{fig:es-inestable}
\end{figure}

\subsection{Dos parámetros: la región de estabilidad de un PI}

Con dos parámetros, las condiciones de la tabla dibujan una \emph{región} en el plano de los parámetros.

\begin{ejemplo}[label=ej:es-pi]{Región de estabilidad de un PI}
En el lazo de la \cref{fig:es-pi-lazo}, $\Gc(s)=\Kc+K_i/s$ y $\Gp(s)=\dfrac{2}{s^3+4s^2+5s+2}=\dfrac{2}{(s+1)^2(s+2)}$. Halle la región del plano $(\Kc,K_i)$ en la que el lazo es estable.
\solucion La ecuación característica es $s(s^3+4s^2+5s+2)+2(\Kc s+K_i)=0$:
\[
  s^4+4s^3+5s^2+(2+2\Kc)\,s+2K_i=0 .
\]
\[
\begin{array}{c|ccc}
  s^4 & 1 & 5 & 2K_i\\
  s^3 & 4 & 2+2\Kc & \\
  s^2 & \dfrac{9-\Kc}{2} & 2K_i\\[6pt]
  s^1 & c_1 & \\
  s^0 & 2K_i
\end{array}
\qquad
c_1=(2+2\Kc)-\frac{16K_i}{9-\Kc}.
\]
Las condiciones son $K_i>0$, $\Kc<9$ y $c_1>0$:
\[
  \alerta{0<K_i<\frac{(1+\Kc)(9-\Kc)}{8},\qquad -1<\Kc<9}.
\]
La frontera es una parábola con máximo $K_i=25/8$ en $\Kc=4$ (\cref{fig:es-region}). Sobre ella, la fila $s^1$ se anula y el auxiliar $\frac{9-\Kc}{2}s^2+2K_i$ da la frecuencia de la oscilación: con $\Kc=5$ y $K_i=3$, $\omega=\sqrt{2\cdot3/2}=1.732$\,rad/s (periodo $3.63$\,s).
\end{ejemplo}

\begin{figure}[ht]
\centering
\includegraphics[width=0.85\linewidth]{cap05/lazo_pi_estab}
\caption{Lazo del ejemplo~\ref{ej:es-pi}, con el PI en forma paralela.}
\label{fig:es-pi-lazo}
\end{figure}

\begin{figure}[ht]
\centering
\includegraphics[width=\linewidth]{cap05/f_region_pi}
\caption{Ejemplo~\ref{ej:es-pi}. Izquierda: región de estabilidad. Derecha: con $\Kc=5$, $K_i=1$ (adentro) es estable, aunque con $\OS=51\%$; $K_i=3$ (frontera) oscila sin amortiguarse; $K_i=4$ (afuera) es inestable.}
\label{fig:es-region}
\end{figure}

\codigorepo{cap05_estabilidad/C5_Ej5_RegionPI.m}{region}{ejemplo~\ref{ej:es-pi}: la región de estabilidad en el plano $(\Kc,K_i)$}

\begin{modelo}[LazoPI.slx]
El script \cmd{crearModelo\_LazoPI.m} arma el lazo de la \cref{fig:es-pi-lazo} con $\Kc$ y $K_i$ como variables del espacio de trabajo. Recorra la región: un punto adentro, uno sobre la frontera y uno afuera, y mida el periodo de la oscilación sostenida.
\end{modelo}

\begin{cuidado}[Estable no quiere decir aceptable]
El punto $(\Kc,K_i)=(5,1)$ está dentro de la región y, sin embargo, la respuesta tiene un sobrepico del $51\%$ y tarda $22$\,s en establecerse. La tabla de Routh dice \emph{si} es estable; para saber \emph{qué tan} estable hace falta algo más.
\end{cuidado}

\subsection{Estabilidad relativa}

La tabla se puede usar también para exigir un margen: que todos los polos tengan parte real menor que $-\sigma$, no solo menor que cero. Basta con desplazar el eje: se reemplaza $s=z-\sigma$ en la ecuación característica y se aplica el criterio al polinomio en $z$. Las raíces con $\Re z>0$ son las que tienen $\Re s>-\sigma$.

\begin{ejemplo}[label=ej:es-relativa]{Polos a la izquierda de $-0.05$}
En el ejemplo~\ref{ej:es-rangokc}, ¿para qué valores de $\Kc$ todos los polos tienen parte real menor que $-0.05$ (constante de tiempo de la envolvente menor que $20$\,s)?
\solucion Con $s=z-0.05$, $900s^3+420s^2+43s+(1+0.8\Kc)$ se convierte en
\[
  900z^3+285z^2+7.75z+(0.8\Kc-0.2125)=0 .
\]
\[
\begin{array}{c|cc}
  z^3 & 900 & 7.75\\
  z^2 & 285 & 0.8\Kc-0.2125\\
  z^1 & \frac{160-48\Kc}{19} & \\
  z^0 & 0.8\Kc-0.2125
\end{array}
\qquad\Longrightarrow\qquad
\alerta{0.266<\Kc<3.33}.
\]
El margen exigido reduce mucho el rango: de $(-1.25,\,23.83)$ a $(0.266,\,3.33)$. En los dos extremos hay un polo exactamente en $-0.05$ (\cref{fig:es-relativa}).
\end{ejemplo}

\begin{figure}[ht]
\centering
\includegraphics[width=0.5\linewidth]{cap05/f_relativa}
\caption{Ejemplo~\ref{ej:es-relativa}: polos del lazo para $\Kc=0.266$ y $\Kc=3.33$ (extremos del rango con margen) y para $\Kc=\Kcr$, con polos sobre el eje imaginario.}
\label{fig:es-relativa}
\end{figure}

\begin{pausa}
Para $\Gma(s)=\dfrac{K}{s(s+1)(s+2)}$ con retroalimentación unitaria, halle el rango de $K$, $\Kcr$ y $\wcr$.
\respuesta{$s^3+3s^2+2s+K=0$. Por \eqref{eq:es-tercer}: $3\cdot2>K$, y $K>0$: $0<K<6$. En $\Kcr=6$ el auxiliar es $3s^2+6=0$: $\wcr=\sqrt2=1.414$\,rad/s.}
\end{pausa}
